\documentclass[12pt]{article}

\usepackage[portuguese]{babel}
\usepackage{float}
\usepackage{listings}
\usepackage{xcolor}
\usepackage{booktabs}
\usepackage[T1]{fontenc}
\usepackage{indentfirst}
\usepackage{amsfonts, amssymb, amsmath}
\usepackage{graphicx}
\usepackage{enumitem}
\usepackage[colorlinks=true, allcolors=blue]{hyperref}
\usepackage[letterpaper,
            top=1.5cm,
            bottom=1.5cm,
            left=2.75cm,
            right=2.75cm,
            marginparwidth=1.75cm]{geometry}
            
\definecolor{azul}{RGB}{0, 90, 160}

\title{Lista 1 -- Exercícios do Capítulo 2}
\author{Prof. Douglas, Guilherme Lopes}
\date{\today}

\begin{document}

\begin{enumerate} [label=\arabic*.]
    \item Usando os dados da Tabela 1, construa a distribuição de freqüências 
    das variáveis:
        \begin{table}[H]
            \centering
            \small
            \setlength{\tabcolsep}{3pt}
            \label{tab:informacoes_empregados}
            \resizebox{\textwidth}{!}{
            \begin{tabular}{c|l|l|c|c|c|c|c}
                \hline
                Nº & Estado civil & Grau de instrução & Nº de filhos &
                Salário & \multicolumn{2}{c}{Idade} & Região de procedência \\
                & & & & (× sal. mín.) & anos & meses & \\
                \hline
                1  & solteiro & ensino fundamental & -- & 4,00  & 26 & 03 & interior \\
                2  & casado   & ensino fundamental & 1  & 4,56  & 32 & 10 & capital \\
                3  & casado   & ensino fundamental & 2  & 5,25  & 36 & 05 & capital \\
                4  & solteiro & ensino médio       & -- & 5,73  & 20 & 10 & outra \\
                5  & solteiro & ensino fundamental & -- & 6,26  & 40 & 07 & outra \\
                6  & casado   & ensino fundamental & 0  & 6,66  & 28 & 00 & interior \\
                7  & solteiro & ensino fundamental & -- & 6,86  & 41 & 00 & interior \\
                8  & solteiro & ensino fundamental & -- & 7,39  & 43 & 04 & capital \\
                9  & casado   & ensino médio       & 1  & 7,59  & 34 & 10 & capital \\
                10 & solteiro & ensino médio       & -- & 7,44  & 23 & 06 & outra \\
                11 & casado   & ensino médio       & 2  & 8,12  & 33 & 06 & interior \\
                12 & solteiro & ensino fundamental & -- & 8,46  & 27 & 11 & capital \\
                13 & solteiro & ensino médio       & -- & 8,74  & 37 & 05 & outra \\
                14 & casado   & ensino fundamental & 3  & 8,95  & 44 & 02 & outra \\
                15 & casado   & ensino médio       & 0  & 9,13  & 30 & 05 & interior \\
                16 & solteiro & ensino médio       & -- & 9,35  & 38 & 08 & outra \\
                17 & casado   & ensino médio       & 1  & 9,77  & 31 & 07 & capital \\
                18 & casado   & ensino fundamental & 2  & 9,80  & 39 & 07 & outra \\
                19 & solteiro & superior            & -- & 10,53 & 25 & 08 & interior \\
                20 & solteiro & ensino médio       & -- & 10,76 & 37 & 04 & interior \\
                21 & casado   & ensino médio       & 1  & 11,06 & 30 & 09 & outra \\
                22 & solteiro & ensino médio       & -- & 11,59 & 34 & 02 & capital \\
                23 & solteiro & ensino fundamental & -- & 12,00 & 41 & 00 & outra \\
                24 & casado   & superior            & 0  & 12,79 & 26 & 01 & outra \\
                25 & casado   & ensino médio       & 2  & 13,23 & 32 & 05 & interior \\
                26 & casado   & ensino médio       & 2  & 13,60 & 35 & 00 & outra \\
                27 & solteiro & ensino fundamental & -- & 13,85 & 46 & 07 & outra \\
                28 & casado   & ensino médio       & 0  & 14,69 & 29 & 08 & interior \\
                29 & casado   & ensino médio       & 5  & 14,71 & 40 & 06 & interior \\
                30 & casado   & ensino médio       & 2  & 15,99 & 35 & 10 & capital \\
                31 & solteiro & superior            & -- & 16,22 & 31 & 05 & outra \\
                32 & casado   & ensino médio       & 1  & 16,61 & 36 & 04 & interior \\
                33 & casado   & superior            & 3  & 17,26 & 43 & 07 & capital \\
                34 & solteiro & superior            & -- & 18,75 & 33 & 07 & capital \\
                35 & casado   & ensino médio       & 2  & 19,40 & 48 & 11 & capital \\
                36 & casado   & superior            & 3  & 23,30 & 42 & 02 & interior \\
                \bottomrule
            \end{tabular}
            }\small
            \caption{Informações sobre estado civil, grau de instrução, número de 
            filhos, salário (expresso como fração do salário mínimo), idade 
            (medida em anos e meses) e procedência de 36 empregados da seção de 
            orçamentos da Companhia MB.}
        \end{table}
        
    \begin{enumerate}[label=\alph*)]
        \item Estado civil.
        \item Região de procedência.
        \item Número de filhos dos empregados casados.
        \item Idade.
    \end{enumerate}
    \vspace{28pt}
    \item Contou-se o número de erros de impressão da primeira página de um jornal 
    durante 50 dias, obtendo-se os resultados abaixo:
    \begin{table}[H]
        \centering
        \begin{tabular}{*{10}{c} }
            8 & 11 & 8 & 12 & 14 & 13 & 11 & 14 & 14 & 15 \\
            6 & 10 & 14 & 19 & 6 & 12 & 7 & 5 & 8 & 8 \\
            10 & 16 & 10 & 12 & 12 & 8 & 11 & 6 & 7 & 12 \\
            7 & 10 & 14 & 5 & 12 & 7 & 9 & 12 & 11 & 9 \\
            14 & 8 & 14 & 8 & 12 & 10 & 12 & 22 & 7 & 15 \\
        \end{tabular}
    \end{table}
    \begin{enumerate}[label=\alph*)]
        \item Represente os dados graficamente
    \end{enumerate}

\end{enumerate}

\end{document}